\documentclass[11pt]{article}
\usepackage{amsmath, amsthm, amssymb}
\usepackage[margin=1.1in]{geometry}
\usepackage{hyperref}

\newtheorem{theorem}{Theorem}
\newtheorem{proposition}[theorem]{Proposition}
\newtheorem{corollary}[theorem]{Corollary}
\newtheorem{lemma}[theorem]{Lemma}
\theoremstyle{definition}
\newtheorem*{conjecture}{Conjecture}
\theoremstyle{remark}
\newtheorem*{remark}{Remark}

\newcommand{\Frob}{\operatorname{Frob}}
\newcommand{\Ind}{\operatorname{Ind}}
\newcommand{\Res}{\operatorname{Res}}
\newcommand{\Hom}{\operatorname{Hom}}
\newcommand{\sgn}{\operatorname{sgn}}
\newcommand{\ch}{\operatorname{ch}}
\newcommand{\triv}{\mathrm{triv}}
\newcommand{\fdeg}{\widetilde{f}}
\newcommand{\qbinom}[2]{\left[{#1 \atop #2}\right]}

\title{The $r$-wreath extension of Theorem D: \\
$S_r$-isotypic decomposition of $R_{(k^r)}$ and\\
a cyclotomic Hilbert-series identity}
\author{Clio}
\date{2026-07-30}

\begin{document}
\maketitle

\begin{abstract}
Fix integers $r, k \ge 1$, let $n = rk$, and let $\alpha = (k^r)$ be the
composition with $r$ parts each equal to $k$. Write
$S_\alpha = S_k^r \subseteq S_n$ for the Young subgroup and
$H = S_k \wr S_r$ for its wreath normaliser. The residual quotient
$H / S_\alpha \cong S_r$ acts on the parabolic coinvariant algebra
$R_\alpha = R_n^{S_\alpha}$ by graded ring automorphisms. We prove:
\begin{enumerate}
\item[\textbf{(A)}] the graded $S_r$-isotypic multiplicities are
$m_\pi(t) = \sum_\lambda \fdeg_\lambda(t) \cdot \langle s_\pi[h_k], s_\lambda\rangle$;
\item[\textbf{(B)}] the graded Frobenius character is
$\Frob_t(R_\alpha) = \sum_\pi m_\pi(t) s_\pi$;
\item[\textbf{(C)}] at $k = 2$, the total Hilbert series admits the
cyclotomic product factorisation
$\binom{2r}{2^r}_t = [r]_{t^2}! \cdot \prod_{i=1}^r [2i-1]_t$.
\end{enumerate}
Result (A)+(B) extends the $r = 2$ Theorem D of \emph{Wreath decomposition
of the collision-regime coset Poincar\'e identity} (2026-07-30) to
arbitrary $r$; result (C) is a clean closed form for the ``collision
regime'' Hilbert series at $k = 2$. We furthermore verify computationally
that at $r = 3, k = 2$ the cyclotomic factorisation of (C) is
\emph{not} aligned with the isotypic decomposition of (A): the factor
$\Phi_3 = [3]_t$ divides $\binom{6}{2,2,2}_t$ but no individual
$m_\pi(t)$.
\end{abstract}

\section{Setup}

Let $S_n$ act on the polynomial ring $\mathbb{C}[x_1,\ldots,x_n]$ by
permutation. Let $I_n = (\mathbb{C}[x_1,\ldots,x_n]_+^{S_n})$ be the
ideal generated by the positive-degree symmetric polynomials, and set
\[
R_n = \mathbb{C}[x_1,\ldots,x_n] / I_n,
\]
the \emph{coinvariant algebra}, a graded $S_n$-module with total Hilbert
series $[n]_t!$. By Chevalley--Shephard--Todd (and Steinberg for the
grading),
\[
R_n \cong \bigoplus_{\lambda \vdash n} \fdeg_\lambda(t) \cdot V_\lambda
\qquad \text{as graded $S_n$-modules,}
\]
where $V_\lambda$ is the irreducible $S_n$-module of shape $\lambda$ and
$\fdeg_\lambda(t)$ is its \emph{fake degree}. In terms of hook lengths,
$\fdeg_\lambda(t) = t^{n(\lambda)} [n]_t! / \prod_{c \in \lambda} [h(c)]_t$
where $n(\lambda) = \sum_i (i-1) \lambda_i$.

Fix a composition $\alpha = (\alpha_1, \ldots, \alpha_r)$ of $n$, and let
$S_\alpha = S_{\alpha_1} \times \cdots \times S_{\alpha_r} \subseteq S_n$
be the corresponding Young subgroup. The \emph{parabolic coinvariant
algebra} is
\[
R_\alpha := R_n^{S_\alpha}.
\]
Its Hilbert series is the $q$-multinomial
\[
\dim_t R_\alpha = \binom{n}{\alpha_1, \ldots, \alpha_r}_t
= \frac{[n]_t!}{[\alpha_1]_t! \cdots [\alpha_r]_t!}.
\]
(Reference: Garsia--Stanton, JLMS 1984; also follows from
$\dim_t R_\alpha = \sum_\lambda \fdeg_\lambda(t) \langle s_\lambda, h_\alpha\rangle$
and the Cauchy identity $\sum_\lambda s_\lambda h_\alpha = h_\alpha$.)

\medskip

Specialise now to $\alpha = (k^r)$, i.e., all parts equal to $k$. The
wreath product
\[
H := S_k \wr S_r = S_\alpha \rtimes S_r
\]
is the normaliser of $S_\alpha$ in $S_{rk}$; the projection
$H \twoheadrightarrow S_r$ has kernel $S_\alpha$.

\section{Main theorem}

\begin{theorem}\label{thm:main}
Let $r, k \ge 1$, $n = rk$, $\alpha = (k^r)$, and $H = S_k \wr S_r$.
\begin{enumerate}
\item[(i)] $H$ preserves $R_\alpha = R_n^{S_\alpha}$ and $S_\alpha$ acts
  trivially, so $H/S_\alpha \cong S_r$ acts on $R_\alpha$ by graded
  ring automorphisms.
\item[(ii)] The graded Frobenius character of $R_\alpha$ as an
  $S_r$-module is
\[
\Frob_t(R_\alpha) \;=\; \sum_{\pi \vdash r} m_\pi(t) \, s_\pi,
\]
with graded isotypic multiplicities
\[
\boxed{\;
m_\pi(t) \;=\; \sum_{\lambda \vdash n} \fdeg_\lambda(t) \cdot
\langle s_\pi[h_k],\, s_\lambda \rangle.
\;}
\]
\item[(iii)] The total Hilbert series is
\[
\dim_t R_\alpha \;=\; \sum_{\pi \vdash r} f^\pi \, m_\pi(t)
\;=\; \binom{n}{k,\,k,\,\ldots,\,k}_t,
\]
where $f^\pi = \dim V_\pi$.
\end{enumerate}
\end{theorem}

\begin{proof}
(i) The subgroup $H \subseteq S_n$ normalises $S_\alpha$, so
$H$ acts on the $S_\alpha$-invariants $R_n^{S_\alpha} = R_\alpha$, and
$S_\alpha$ acts trivially on its own invariants. Consequently the
action factors through $H/S_\alpha \cong S_r$.

\medskip

(ii) Applying the exact functor $(-)^{S_\alpha}$ (since we are in
characteristic zero) to the CST decomposition,
\[
R_\alpha \;=\; \bigoplus_{\lambda \vdash n} \fdeg_\lambda(t) \cdot
V_\lambda^{S_\alpha}
\qquad \text{as graded $S_r$-modules,}
\]
where the $S_r$-action on each $V_\lambda^{S_\alpha}$ comes from the
$H$-action on $V_\lambda|_H$ (well-defined by the same normalisation
argument).

Fix $\pi \vdash r$. Let $\widetilde{V}_\pi := V_\pi \boxtimes \triv^{\otimes r}$
denote the $H$-module obtained by pulling back $V_\pi$ along
$H \twoheadrightarrow S_r$; equivalently, $S_\alpha$ acts trivially and
$S_r$ acts by $V_\pi$. Then:
\begin{align*}
\Hom_{S_r}(V_\pi, V_\lambda^{S_\alpha})
&= \Hom_H(\widetilde{V}_\pi, V_\lambda^{S_\alpha}) \tag{$S_\alpha$ trivial on both}\\
&= \Hom_H(\widetilde{V}_\pi, V_\lambda|_H) \tag{image is $S_\alpha$-fixed}\\
&= \Hom_{S_n}(\Ind_H^{S_n}(\widetilde{V}_\pi),\, V_\lambda). \tag{Frobenius reciprocity}
\end{align*}
The Frobenius character of the induced $S_n$-module
$\Ind_H^{S_n}(\widetilde{V}_\pi)$ is the plethysm $s_\pi[h_k]$: this is
the classical formula for wreath induction of a representation trivial
on the base subgroup (Macdonald, \emph{Symmetric Functions and Hall
Polynomials}, Ch.~I, App.~A, eq. (A.11); Stanley, \emph{EC2}, ex.~7.75).
Hence
\[
\dim \Hom_{S_r}(V_\pi, V_\lambda^{S_\alpha})
= \langle s_\pi[h_k],\, s_\lambda \rangle.
\]
Summing $\lambda$ with graded weight $\fdeg_\lambda(t)$:
\[
m_\pi(t) := \dim_t \Hom_{S_r}(V_\pi, R_\alpha)
= \sum_\lambda \fdeg_\lambda(t) \, \langle s_\pi[h_k], s_\lambda\rangle.
\]

\medskip

(iii) Direct consequence of $R_\alpha \cong \bigoplus_\pi m_\pi(t) \cdot V_\pi$
and $\dim V_\pi = f^\pi$; equality with the $q$-multinomial is the
classical Hilbert-series formula for parabolic coinvariants recalled in
\S 1.
\end{proof}

\section{Cyclotomic factorisation at $k = 2$}

\begin{theorem}\label{thm:cyclotomic}
For every $r \ge 1$,
\[
\binom{2r}{2, 2, \ldots, 2}_t \;=\; [r]_{t^2}! \cdot \prod_{i=1}^{r} [2i-1]_t.
\]
\end{theorem}

\begin{proof}
Since $[1]_t! = [2]_t$ (with $[2]_t = 1+t$), we have $[2]_t! = [2]_t$
and $([2]_t!)^r = [2]_t^r$. Thus
\[
\binom{2r}{2^r}_t = \frac{[2r]_t!}{[2]_t^r}.
\]
Split $[2r]_t! = \prod_{m=1}^{2r} [m]_t$ into even and odd factors:
\[
[2r]_t! \;=\; \Big(\prod_{i=1}^{r} [2i]_t\Big) \cdot \Big(\prod_{i=1}^{r} [2i-1]_t\Big).
\]
The elementary identity
\[
[2i]_t \;=\; [2]_t \cdot [i]_{t^2}
\qquad
\left(\text{since } \frac{1-t^{2i}}{1-t} = \frac{1-t^2}{1-t} \cdot \frac{1-t^{2i}}{1-t^2}\right)
\]
gives $\prod_{i=1}^r [2i]_t = [2]_t^r \cdot [r]_{t^2}!$. Substituting
and cancelling $[2]_t^r$:
\[
\binom{2r}{2^r}_t \;=\; [r]_{t^2}! \cdot \prod_{i=1}^{r} [2i-1]_t. \qedhere
\]
\end{proof}

\begin{remark}[Combinatorial interpretation]
The specialisation $t = 1$ recovers $\binom{2r}{2^r} = r! \cdot (2r-1)!!$,
the standard split of ordered pair-decompositions of $\{1,\ldots,2r\}$
into (matching selection) $\times$ (pair ordering). The factor
$[r]_{t^2}!$ is the pair-ordering contribution in $t^2$: swapping two
adjacent pairs contributes $2$ to the major index (each pair is stored
increasingly, so a swap crosses each side of the other pair once).
\end{remark}

\begin{remark}[Cyclotomic vs.\ isotypic decomposition]
The factorisation of Theorem~\ref{thm:cyclotomic} is a \emph{Hilbert-series}
identity and does not, in general, align with the $S_r$-isotypic
decomposition of Theorem~\ref{thm:main}. For $(r, k) = (3, 2)$ we verify
computationally (\S 4) that $[3]_t = \Phi_3$ divides
$\binom{6}{2,2,2}_t = [5]_t [3]_t [3]_{t^2} [2]_{t^2}$ but does
\emph{not} divide any individual $m_\pi(t)$. In contrast, the factor
$[5]_t = [2r-1]_t$ turns out to divide \emph{every} $m_\pi(t)$; this
divisibility appears empirically to hold whenever $k = 2$ (verified
$r = 2, 3, 4$) and additionally at some higher $(r, k)$ (e.g.\ $(3,3)$
but not $(2,3)$), but does not follow from the framework of this note.
\end{remark}

\section{Verification at $(r, k) = (3, 2)$}

The plethysms $s_\pi[h_2]$ for $\pi \vdash 3$ are (computed in
\texttt{probes/2026-07-30-r3-wreath-extension/plethysm.py}):
\begin{align*}
s_{(3)}[h_2] &= s_{(6)} + s_{(4,2)} + s_{(2,2,2)}, \\
s_{(2,1)}[h_2] &= s_{(5,1)} + s_{(4,2)} + s_{(3,2,1)}, \\
s_{(1,1,1)}[h_2] &= s_{(4,1,1)} + s_{(3,3)}.
\end{align*}
The first is the classical Foulkes formula for $h_3[h_2]$. Summed with
$S_3$-dimension weights ($f^{(3)} = 1$, $f^{(2,1)} = 2$, $f^{(1^3)} = 1$):
\[
s_{(3)}[h_2] + 2 s_{(2,1)}[h_2] + s_{(1,1,1)}[h_2]
= s_{(6)} + 2 s_{(5,1)} + 3 s_{(4,2)} + s_{(4,1,1)} + s_{(3,3)} + 2 s_{(3,2,1)} + s_{(2,2,2)} = h_2^3,
\]
consistent with $\Frob(\Ind_{S_\alpha}^{S_6}(\triv)) = h_2^3$.

The fake degrees of the six relevant shapes (via $\fdeg_\lambda(t) = t^{n(\lambda)}[6]_t!/\prod [h(c)]_t$):
\[
\begin{array}{r|l}
\lambda & \fdeg_\lambda(t) \\ \hline
(6)     & 1 \\
(5,1)   & t + t^2 + t^3 + t^4 + t^5 \\
(4,2)   & t^2 + t^3 + 2t^4 + t^5 + 2t^6 + t^7 + t^8 \\
(4,1,1) & t^3 + t^4 + 2t^5 + 2t^6 + 2t^7 + t^8 + t^9 \\
(3,3)   & t^3 + t^5 + t^6 + t^7 + t^9 \\
(3,2,1) & t^4 + 2t^5 + 2t^6 + 3t^7 + 3t^8 + 2t^9 + 2t^{10} + t^{11} \\
(2,2,2) & t^6 + t^8 + t^9 + t^{10} + t^{12}
\end{array}
\]

Applying Theorem~\ref{thm:main}(ii):
\begin{align*}
m_{(3)}(t) &= \fdeg_{(6)} + \fdeg_{(4,2)} + \fdeg_{(2,2,2)} \\
&= 1 + t^2 + t^3 + 2t^4 + t^5 + 3t^6 + t^7 + 2t^8 + t^9 + t^{10} + t^{12},\\
m_{(2,1)}(t) &= \fdeg_{(5,1)} + \fdeg_{(4,2)} + \fdeg_{(3,2,1)} \\
&= t + 2t^2 + 2t^3 + 4t^4 + 4t^5 + 4t^6 + 4t^7 + 4t^8 + 2t^9 + 2t^{10} + t^{11},\\
m_{(1,1,1)}(t) &= \fdeg_{(4,1,1)} + \fdeg_{(3,3)} \\
&= 2t^3 + t^4 + 3t^5 + 3t^6 + 3t^7 + t^8 + 2t^9.
\end{align*}
Weighted sum:
\begin{align*}
m_{(3)}(t) + 2 m_{(2,1)}(t) + m_{(1,1,1)}(t)
&= 1 + 2t + 5t^2 + 7t^3 + 11t^4 + 12t^5 + 14t^6 \\
&\quad {}+ 12t^7 + 11t^8 + 7t^9 + 5t^{10} + 2t^{11} + t^{12} \\
&= \binom{6}{2,2,2}_t \\
&= [5]_t \cdot [3]_t \cdot [3]_{t^2} \cdot [2]_{t^2}. \qquad \checkmark
\end{align*}

Palindromicity of each $m_\pi(t)$ is inherited from the palindromicity
of every $\fdeg_\lambda(t)$ (the coinvariant algebra is a Frobenius
algebra with top degree $\binom{n}{2}$; for $n = 6$, top degree $15$).
The individual multiplicities are centred at degrees
$6$, $6$, $6$ respectively (their supports are $[0,12]$, $[1,11]$,
$[3,9]$).

\section{Computational verification}

The identity of Theorem~\ref{thm:main}(iii),
\[
\sum_{\pi \vdash r} f^\pi \, m_\pi(t) = \binom{rk}{k^r}_t,
\]
has been verified in \texttt{probes/2026-07-30-r3-wreath-extension/verify\_full.py}
for
\[
(r, k) \in \{(2,2),\, (2,3),\, (3,2),\, (2,4),\, (3,3),\, (4,2)\}.
\]
The cyclotomic identity of Theorem~\ref{thm:cyclotomic} was verified
symbolically for $r = 1, 2, \ldots, 6$ in the same script.

\section{Remarks and follow-up questions}

\begin{remark}[Relation to Levic\'an--Romero]
The framework of Levic\'an--Romero (arXiv:2504.19008) develops
$\mathbb{Z}_k \wr S_n$-class functions for wreath Macdonald theory.
The present setup ($S_k \wr S_r$ acting on $R_{rk}^{S_k^r}$) is
classical rather than wreath-Macdonald, but the underlying combinatorics
(plethysm $s_\pi[h_k]$, Frobenius reciprocity for $H$) is a
degenerate specialisation of their machinery. A bigraded $(q,t)$-lift
following Szendr\H oi arXiv:2602.15017 or Romero--Wen
arXiv:2505.01732 is the natural next question.
\end{remark}

\begin{remark}[$[2r-1]_t$-divisibility of $m_\pi(t)$]
Empirically, $[2r-1]_t$ divides every $m_\pi(t)$ when $k = 2$
(verified $r = 2, 3, 4$), and also at $(r, k) = (3, 3)$, but does
not at $(2, 3)$ or $(2, 4)$. This suggests a rank-two phenomenon
governed by the divisibility of the fake degrees $\fdeg_\lambda(t)$ by
cyclotomic factors indexed by hook data. A conceptual explanation
(via, e.g., Springer-fibre cohomology or Solomon's descent-algebra
picture) is not attempted here.
\end{remark}

\begin{remark}[Sprint context]
This note completes Theorem D of the byproduct paper
\emph{Modified Hall--Littlewood boundary} (\S 5.2) at general $r$.
The $r = 2$ case was proved in the companion note
\emph{Wreath decomposition of the collision-regime coset Poincar\'e
identity} (2026-07-30, morning); this extension supplies the
general-$r$ version of Conjectures A, B, and (partially) C from
\texttt{PROVE.md} of 2026-07-30 late-wake. Conjecture C is proved
as a Hilbert-series identity but shown to be \emph{not} the
$S_r$-isotypic decomposition.
\end{remark}

\end{document}
